% ============================================================================
% Fragment dành cho chương Kết quả. Chèn sau bảng kết quả định lượng tổng quan.
% Preamble cần: booktabs, tabularx, array.
% ============================================================================

\subsection{Chi phí tính toán và hành vi huấn luyện}

\subsubsection{Chi phí huấn luyện thực đo}

Bảng~\ref{tab:wallclock} tổng hợp thời gian huấn luyện quan sát được trên cùng
một RTX 3060. Các giá trị không bao gồm thời gian tiền xử lý, sinh dự đoán hoặc
chấm điểm.

\begin{table}[tbp]
  \centering
  \small
  \renewcommand{\arraystretch}{1.18}
  \caption{Tổng thời gian huấn luyện ghi nhận của ba hệ thống.}
  \label{tab:wallclock}
  \begin{tabularx}{\textwidth}{@{}>{\bfseries}p{3.4cm}p{3.1cm}X@{}}
    \toprule
    Hệ thống & Tổng wall-clock & Ghi chú \\
    \midrule
    E1
      & $\approx$ 3.209,5 s ($\approx$ 53,5 phút)
      & Cộng khoảng 90 s của lần chạy bị ngắt và 3.119,5 s của đoạn chạy tiếp
        từ checkpoint tại epoch 5 / update 579 \\
    E3-custom
      & 4.758,7 s ($\approx$ 79,3 phút)
      & Một lượt chạy; chạm \texttt{max\_update=50000} \\
    E2
      & 33.940,7 s ($\approx$ 9,43 giờ)
      & Cộng đoạn đầu 20.328,3 s và đoạn chạy tiếp 13.612,4 s từ checkpoint
        546 đến step 910 \\
    \bottomrule
  \end{tabularx}
\end{table}

Tổng của E1 là tổng thời gian còn quan sát được qua hai đoạn, không phải phép đo
của một lượt huấn luyện sạch từ đầu đến cuối. Tương tự, tổng E2 được tái dựng
từ thời gian đã ghi cho đoạn đầu và \texttt{train\_runtime} của đoạn chạy tiếp;
log chi tiết của đoạn đầu không còn trong artifact hiện tại. Vì vậy, các con số
này phù hợp để mô tả chi phí đã ghi nhận, nhưng không nên được diễn giải như
benchmark tốc độ chuẩn hoá giữa các kiến trúc.

\subsubsection{Ảnh hưởng quan sát được của giới hạn bộ nhớ}

Khi E3-custom bổ sung danh sách gợi ý Hán tự, số token nguồn tăng xấp xỉ
3,9 lần. Dưới cùng giới hạn \texttt{max\_tokens\_per\_gpu=1000}, số câu trung
bình trong mỗi update giảm từ 132,5 ở E1 xuống 33,3 ở E3-custom. Vì vậy, E1 và
E3-custom không chỉ khác nhau ở nội dung đầu vào: số câu đóng góp vào mỗi lần
cập nhật cũng khác do giới hạn bộ nhớ. Đây là một confound của phép so sánh cấp
hệ thống, không phải bằng chứng riêng về hiệu quả của cơ chế truy hồi.

\subsubsection{Diễn biến validation và lựa chọn checkpoint của E2}
\label{sec:e2-budget}

E2 ban đầu được huấn luyện trong 3 epoch, tương ứng 546 optimizer step, rồi chạy
tiếp đến 5 epoch và 910 step từ checkpoint 546. Với
\texttt{eval\_steps=200}, toàn bộ quá trình chỉ có bốn điểm đánh giá validation
(Bảng~\ref{tab:e2-budget}).

\begin{table}[tbp]
  \centering
  \small
  \renewcommand{\arraystretch}{1.18}
  \caption{Các điểm đánh giá validation của E2 qua 910 step.}
  \label{tab:e2-budget}
  \begin{tabularx}{\textwidth}{@{}rrrX@{}}
    \toprule
    \textbf{Step} & \textbf{Epoch} & \textbf{\texttt{eval\_loss}}
      & \textbf{Giai đoạn} \\
    \midrule
    200 & 1,10 & 1,4444 & Ba epoch ban đầu \\
    400 & 2,20 & \textbf{1,4322} & Ba epoch ban đầu; checkpoint được chọn \\
    600 & 3,30 & 1,4659 & Đoạn chạy tiếp \\
    800 & 4,40 & 1,5815 & Đoạn chạy tiếp \\
    \bottomrule
  \end{tabularx}
\end{table}

Trong bốn điểm được đo, \texttt{eval\_loss} thấp nhất ở step 400 và cao hơn ở
step 600 và 800. Do Trainer dùng \texttt{load\_best\_model\_at\_end} theo
\texttt{eval\_loss}, adapter xuất ra sau step 910 vẫn mang trọng số của
checkpoint 400. SHA-256 của \texttt{adapter\_model.safetensors} xuất ra trùng
với file tương ứng trong checkpoint 400, trong khi khác checkpoint 800 và 910.

Kết quả này cho thấy phần ngân sách kéo dài không thay đổi checkpoint được chọn;
nó không tạo thành một hệ thống E2 mới. Tuy nhiên, bốn điểm đánh giá là một lưới
thưa và tiêu chí lựa chọn là loss, không phải BLEU. Vì vậy, bằng chứng chỉ hỗ trợ
nhận định rằng validation loss đo được tăng sau step 400; nó không xác định
chính xác thời điểm bắt đầu overfitting và cũng không chứng minh bản dịch ở 5
epoch kém hơn theo BLEU.
